\documentclass[10pt,a4paper]{article}
\usepackage[T1]{fontenc}
\usepackage{lmodern,amsmath,amssymb,booktabs,microtype}
\usepackage[margin=24mm]{geometry}
\usepackage[hidelinks,hypertexnames=false]{hyperref}
\pdfinfoomitdate=1
\pdftrailerid{}
\pdfsuppressptexinfo=15
\setlength{\parindent}{0pt}
\setlength{\parskip}{5pt}
\begin{document}
\begin{center}
{\Large Hydrogen line escape and the shared radiation field}\\[4pt]
A separate internal-transfer closure for the coupled composition model
\end{center}

\section{What is solved}
\path{escape.hydrogen_escape_probability.HydrogenEscapeProbabilityModel}
accepts temperature $T$, hydrogen level populations $n_i$, proton density
$n_p$, and electron density $n_e$ independently in every slab cell. These
quantities may be unknowns of a larger composition and thermal root.
The class evaluates hydrogen's contributions to that root; its own
\texttt{solve} method closes only a fixed-temperature, pure-hydrogen slab.
No Cloudy process or stored Cloudy population is used at runtime.

The selectable \texttt{anderson\_lebedev} collision model uses Anderson
data for $n=4,5$ and C23's Lebedev--Beigman shell formula~\cite{lebedev} for
$n>5$. A collapsed upper shell feeds a resolved lower orbital state with
fraction $(2l+1)/n^2$, rather than its radiative branching fraction. C23's
electron--atom reduced-mass correction is included for these added-shell
rates, and upward rates retain detailed balance at the actual atom energies.
With \texttt{resolved\_nmax=4} through 10, the original Anderson orbital rows
are retained, including radiatively forbidden electron transitions.
Only collapsed states sum these rows. The same density-dependent proton
mixing kernel connects adjacent orbital states within each resolved shell.
The existing \texttt{anderson\_n45} choice permits resolution only through 5.
Above 5, only
\texttt{anderson\_lebedev} is permitted. Its shell collision strength is
allocated over \emph{both} resolved states using $(2l+1)/n^2$ for each,
and every distinct-$n$ pair is included even if its radiative dipole is zero.
This retains C23's shell approximation and the actual native energy gaps
for inverse detailed balance. The legacy \texttt{vs80} choice retains
resolution through $n=3$ only. Proton orbital mixing is unchanged;
He-ion mixing and heavy-particle $n$-changing collisions are separate omissions.
This removes a measured high-shell rate discrepancy; it does not establish
convergence of the finite atom or validate all low-level data.

The optional \texttt{topoff=True} closure transfers C23's omitted-shell
recombination to the highest collapsed shell. Its temperature-dependent
factor multiplies every bound--free forward/inverse and thermal path,
including signed shared-line photons and OTS absorption. The highest
collisional-ionization and inverse three-body coefficients are both
enhanced by 100; that artificial pair contributes no kinetic heat.
The solar-model supplement defines the factor, domain and prescribed-LTE
mask. This fixed-atom option rejects pressure-lowered active cells and is
restricted to 4000--14000 K; it does not establish shell convergence.

The model is an explicit alternative to the existing full-$J$ CRD transfer.
Line photons that repeatedly scatter are represented by escape and destruction
probabilities, following adapted Cloudy C23.01 operators~\cite{cloudy}.
They must not also enter the excitation rate as a trapped local intensity.
The observer remains separate: changing this internal closure does not change
its line-profile width convention or its integration over viewing rays.

\section{The radiation--population loop}
At a transition $u\leftrightarrow l$, the line opacity depends on the current
populations,
\begin{equation}
 \kappa_{ul}=\frac{c^2A_{ul}}{8\pi\nu_{ul}^2\sqrt\pi\,\Delta\nu_D}
 \left(\frac{g_u}{g_l}n_l-n_u\right),\qquad
 \Delta\nu_D=\frac{\nu_{ul}}c\sqrt{\frac{2k_{\rm B}T}{m_{\rm H}}}.
 \label{eq:opacity}
\end{equation}
Columns are sums of these opacities across the actual cells, including
fine-component overlap. The optical depth toward the illuminated surface is
an \emph{exclusive} sum: the illuminated face has exactly zero overlying
column. Pump shielding separately includes a finite-cell correction.
The finite-slab default uses $\tau_{\rm tot}=\sum_c\kappa_c\Delta z_c$.
The explicit \texttt{line\_optical\_depth\_scale=2} option doubles that total,
matching sunbather's declared Cloudy \texttt{double optical depths} prescription.
It represents an additional column beyond the computed base; it is not the
open-back boundary or a parameter fitted to the spectrum.

The explicit \texttt{line\_sampling="average"} alternative evaluates the
operators inside each cell. Gauss nodes $u_q\in(0,1)$ and weights
$\sum_qw_q=1$ give
\begin{equation}
 \tau_q=\tau_{\rm front}+u_q\Delta\tau,\qquad
 \overline P=\sum_q w_q P(\tau_q,\tau_{\rm tot}),\qquad
 \overline{S\eta}=\sum_qw_q S(\tau_q)\eta_{\rm cont}(u_q).
 \label{eq:cell_average}
\end{equation}
Escape, destruction, scattering, stimulated rates and returned OTS photons
use the same weights. Stellar attenuation and the one-way diffuse formal
solution supply the continuum occupation at each node. Point shielding is
used there: multiplying separately averaged intensity and shielding, or
applying C23's finite-cell shielding correction again, would be inconsistent.
The \texttt{front} option retains the preceding C23-style discretization.
Both must approach the same thin-cell limit. The gain-domain check includes
cell endpoints; averaging does not relax it. This option changes spatial
sampling, not the continuum closure or redistribution approximation.

\paragraph{Resolving a thin illuminated skin.}
Uniform Gauss nodes can miss the thin region where an optically thick cell
absorbs stellar UV photons. The optional
\texttt{line\_quadrature="optical\_depth"} places more nodes near both faces
and splits the cell at the shielding joins and the Ly$\alpha$ continuum
switch $\beta\sqrt\pi\tau=10^{-3}$. For a segment $[a,a+d]$, map ordinary
Gauss nodes $x_q\in(0,1)$, with weights $w_q$, to its two halves:
\begin{align}
 L&=\log(1+S),& g(x)&=\frac{e^{xL}-1}{S},\qquad S>0,\nonumber\\
 u_q^-&=a+\frac d2g(x_q),&
 u_q^+&=a+d\left[1-\frac12g(x_q)\right],\nonumber\\
 \widetilde w_q^\pm&=\frac{dw_qL e^{x_qL}}{2S},&
 W_q^\pm&=\frac{\widetilde w_q^\pm}
                  {\sum_{\text{segments},q,\pm}\widetilde w_q^\pm}.
 \label{eq:optical_depth_quadrature}
\end{align}
The positive, smooth scale $S$ depends on the segment's Ly$\alpha$ depth,
its overlying columns and the continuum depths actually sampled by the
hydrogen lines. The scale changes node placement only. All original point
operators and opacity checks are retained. The returned raw weight-sum
error makes loss of constant-function accuracy visible before normalization.
The order is the number of nodes \emph{per half-segment}, not per cell.

The same cellwise nodes and weights determine hydrogen rates, returned
photons and direct Ly$\alpha$--Fe exchange.
\texttt{hydrogen\_line\_quadrature\_order} allows an independent hydrogen
refinement. Ordinary metal lines default to their separate Gauss rule;
\texttt{metal\_line\_quadrature="optical\_depth"} instead maps each line's
own columns. It splits both faces at the PRD join $a\tau=1$ or the applicable
Doppler-fit joins, and at zero column for an allowed inverted cell. Signed
cell increments retain their sign in the split positions. The original
point kernels and gain guards are unchanged. A weighted scan accumulates
the rates and photon returns without storing every node's Fe-line arrays.
The H--O I overlap solve still replaces
the isolated Ly$\beta$ operator. A successful quadrature check does not
establish convergence of the complete atmosphere or of ordinary metal lines.

The metal rule evaluates the actual diffuse field at each line's moving
node. For each of the two bracketing continuum frequencies $k$ and direction
$\mu_a$, first compute
\begin{equation}
 I_{ak}(u)=I_{ak}^{\rm in}e^{-u\tau_{ak}}
       +uE_{ak}\varphi(u\tau_{ak}),\qquad
 \varphi(x)=\frac{1-e^{-x}}x,\quad \varphi(0)=1,
 \label{eq:line_node_continuum}
\end{equation}
where $\tau_{ak}=\chi_k\Delta z/\mu_a$ and
$E_{ak}=\eta_k\Delta z/\mu_a$. Only then interpolate in frequency and
sum the angular weights. The separately sampled stellar beam uses its
exact line opacity. This preserves the existing one-way diffuse closure;
it does not replace it with bidirectional or moving-medium line transport.

Splitting at the moving Ly$\alpha$ switch also matters for derivatives.
For a parameter $\theta$ and moving interior boundaries $u_b(\theta)$,
\begin{equation}
 \frac{d}{d\theta}\int_0^1 Q(u,\theta)\,du
 =\sum_s\int_{u_s}^{u_{s+1}}\partial_\theta Q\,du
  +\sum_b\left[Q(u_b^-,\theta)-Q(u_b^+,\theta)\right]
                \frac{du_b}{d\theta}.
 \label{eq:moving_quadrature_boundary}
\end{equation}
Automatic differentiation includes the split endpoints, mapped nodes and
weights. Differentiating only fixed nodes would miss the second term when
the source prescription has a jump. This treatment does not smooth that
prescription or claim differentiability where a split changes topology.

The effective radiative rates are
\begin{align}
 p_{lu}&=A_{ul}\frac{g_u}{g_l}\,S(\tau,a)\,
                 \eta_{\rm incident+forward\ diffuse},\\
 R_{lu}^{\rm rad}&=p_{lu},\qquad
 R_{ul}^{\rm rad}=A_{ul}(P_{\rm esc}+P_{\rm dest}+P_e)
                         +\frac{g_l}{g_u}p_{lu}.
 \label{eq:rates}
\end{align}
$S$ is the corrected Federman shielding function, $P_e$ is escape by electron
scattering, and $P_{\rm dest}$ is absorption by the continuum. Ly$\alpha$ uses
the special Hummer--Kunasz prescription; other Lyman lines use the Doppler-core
fit, and subordinate hydrogen lines use the complete-redistribution wing fit.
The separate \path{line_escape_probability} module also exposes the PRD fit
needed by metal lines. Internal damping uses the upper-state lifetime, as in
C23; the observer retains both upper and lower lifetimes.

\paragraph{A known infrared limit of the electron-escape closure.}
The inherited C23 prescription uses
\begin{equation}
 P_e=\frac{\kappa_e}{\kappa_e+\kappa_*}
           \max(0,1-P_{\rm esc}-P_{\rm dest}),\qquad
 \kappa_*=\begin{cases}
 \kappa_{\rm net},&\kappa_{\rm net}>0,\\
 \kappa_{\rm lower},&\kappa_{\rm net}\leq0.
 \end{cases}
 \label{eq:electron_escape_branch}
\end{equation}
Here $\kappa_{\rm net}$ is the multiplet opacity including stimulated
emission, whereas $\kappa_{\rm lower}$ includes absorption only. These
branches need not meet at zero net opacity. A high-shell infrared transition
in the eight-cell continuation crosses this boundary and produces a finite
rate jump. Agreement with the inherited implementation is therefore not
evidence that this closure is adequate there.

The original line--continuum escape treatment~\cite{netzer} uses net line
opacity, a total line-plus-continuum escape probability and a negligible
continuum source. Those assumptions do not justify replacing net opacity by
absorption-only opacity throughout an infrared problem. A candidate repair
must include the actual continuum field and scattering source, preserve
photon and energy accounting, and recover the finite net radiative loss
$n_u A_{ul}$ when $n_lB_{lu}-n_uB_{ul}=0$. This is an unresolved model
obligation, not a claim that the present escape-probability closure meets
those limits.

For a finite positive cell depth, C23's shielding average uses
$d=\Delta\tau/(1+\tau)$ and a factor $\log(1+d)/d$.
When $d<0$, its explicit \texttt{SMALLFLOAT} branch evaluates to zero in
the reference double-precision implementation. The native operator preserves
that branch; evaluating \texttt{log1p} on the negative $d$ instead can fail
near $\tau=-1$. At exactly $\tau=-1$ the reference is undefined; the native
operator uses the continuous zero extension from neighboring finite-cell
values. This is a declared numerical extension, not a smooth-derivative or
physical-maser claim. The independent gain guard at $\tau\leq-30$ remains.

\paragraph{An explicit shielding-fit control.}
The source's positive core uses $e^{-b\tau}$ for $\tau<2$, then
$c_1\tau^{-p_1}$, $c_2\tau^{-p_2}$ and $c_3\tau^{-p_3}$ across
joins at 2, 10 and 100. Its fitted coefficients leave finite jumps, which
can make the coupled residual jump across zero without a root at the join.
The optional \texttt{shielding\_model="continuous\_core"} instead defines
\begin{equation}
 c_1=e^{-2b}2^{p_1},\qquad
 c_2=c_1 10^{p_2-p_1},\qquad c_3=c_2 100^{p_3-p_2},
 \label{eq:connected_core}
\end{equation}
with the original exponents, $b=0.66666$, and unchanged wings. The positive
core changes by at most 1.732\%. This bounds only the operator change, not
the resulting thermal or spectral error. The original \texttt{cloudy}
option remains the default. These three joins become continuous but retain
kinks; the separate negative-to-positive branch jump at $\tau=0$ also
remains. This control is not a globally smooth transfer model.

\paragraph{Continuous Doppler escape control.}
The independent \texttt{core\_escape\_model="continuous"} approximates the
angle-averaged Doppler integral
\begin{equation}
 \beta_D(\tau)=\frac{2}{\sqrt\pi}\int_0^\infty e^{-x^2}
 E_2(\tau e^{-x^2})\,dx,\qquad
 E_2(z)=\int_0^1 e^{-z/\mu}\,d\mu.
 \label{eq:doppler_escape}
\end{equation}
It retains the $\tau\log\tau$ thin limit instead of C23's discontinuous
linear shortcut at $\sqrt\pi\tau=0.01$. A cubic smoothstep joins the two
remaining fits across $10<\sqrt\pi\tau<12$. The sampled independent
integrals differ by at most 0.653\%; this does not improve the separate
wing prescription's stated 20\% uncertainty. The original \texttt{cloudy}
choice remains the default. The value at zero is one, but the positive-side
derivative has a physical logarithmic singularity. A constant zero column
has a zero tangent in code; global differentiability is not claimed.

Radiation changes the populations, which change opacity and hence radiation.
Every evaluation recomputes this loop. A root satisfies statistical equilibrium,
\begin{equation}
 0=\sum_j(n_jR_{ji}-n_iR_{ij})+n_en_p\alpha_i+n_e^2n_p\beta_i
       -n_i\Gamma_i-n_en_iq_{i,\rm ion}.
 \label{eq:se}
\end{equation}
Proton orbital-mixing coefficients are multiplied by $n_p$, even when metal
electrons make $n_e\ne n_p$. Implicit differentiation applies to the complete
selected residual: $d\boldsymbol{x}/d\theta=-F_x^{-1}F_\theta$.
The shared \path{escape.nonlinear} module offers Newton--GMRES using exact
JAX Jacobian--vector products. Its linear tolerance controls only a trial
step; the nonlinear residual has its own acceptance criterion. Tangent and
transpose solves check their residuals inside the linear-solve primitive,
so reverse-mode differentiation does not pass through Newton iterations.

\section{Destroyed photons go to their actual absorber}
Line destruction produces
\begin{equation}
 D_{ul}=n_uA_{ul}P_{\rm dest},\qquad
 \Phi_{ul}^{\rm OTS}=\frac{D_{ul}}{\kappa_{\rm abs}(\nu_{ul})},\qquad
 \Gamma_{i,ul}^{\rm OTS}=\sigma_i(\nu_{ul})\Phi_{ul}^{\rm OTS}.
 \label{eq:ots}
\end{equation}
Here $D$ has units photons cm$^{-3}$ s$^{-1}$ and $\Phi$ photons
cm$^{-2}$ s$^{-1}$ integrated over the line. The common denominator includes
hydrogen, metals and molecules. The supplying network evaluates each species'
opacity at the same frequency and receives its share. Stimulated recombination
uses the same field and Milne inverse coefficient. Photoelectron heating uses
$h\nu-\chi_i$, not the full destroyed-photon energy $h\nu$.
Scattering removal is not thermal absorption.
\texttt{photoionization\_from\_photon\_flux} supplies the matching hydrogen
rates and kinetic ledger for photons returned by another species, ensuring
that metal-to-hydrogen feedback updates both populations and heat.

The returned hydrogen kinetic ledger includes collisional excitation and
ionization, three-body recombination, continuum photoelectrons and radiative
recombination, plus the hydrogen share of OTS absorption. A shared free-free
provider can select \texttt{include\_free\_free=False}; no LTE H-minus opacity
or cooling is inserted by this class.

\section{Continuum, boundary data and limits}
The current continuum implementation uses the existing formal slab solver.
Hydrogen bound-free gain is suppressed per level, as in C23. Its optional
scattering attenuation multiplies transmission through a cell by
$(1+\kappa_{\rm sct}\Delta z)^{-1}$; this is the C23 attenuation approximation,
not conservative angular scattering transport. Consequently this combination
must be compared with C23's default outward/OTS recombination-continuum
closure rather than assumed identical to it.

Spontaneous $2s$ two-photon emission is transported using the C23 natural cubic
spline of Goldman--Drake's nine hydrogen shape samples~\cite{goldman}.
Its quadrature produces exactly two photons per decay; the integrated photon
energy is checked independently. Induced two-photon processes are disabled,
as in C23's default settings.

A normal stellar beam is specified by $J_\nu$; the dark-face boundary is an
incoming isotropic $I_\nu$. For narrow stellar absorption features, supply
\texttt{incident\_line\_mean\_intensity} directly at the line frequencies.
Interpolation from a sparse continuum grid must not erase those features.
Additional continuum thresholds split the frequency quadrature explicitly.

Primitive parity, a converged population residual, and a finite derivative
are separate checks. They do not establish spatial convergence, agreement
with a solar-abundance atmosphere, or the adequacy of a lower-atmosphere join.
In particular, uniformly thick surface cells cannot stand in for a resolved
illuminated boundary. Controlled external references and failed comparisons
are retained separately from runtime model data.
An inverted coarse cell can also lack feedback from its own negative column
under the exclusive-sum closure. Spatial refinement must resolve that
feedback; increasing the gain guard or dropping the offending transition
does not establish a solution.

\section{A conservative H--O I overlap group}
\texttt{LineOverlapGroup} replaces the isolated Ly$\beta$ and six nearby
O I operators with a common, frequency-resolved, static CRD transfer solve.
The actual line centers and thermal masses are retained. This is a different
closure from C23's CORE/PRD fits; in particular, zero O I recovers this same
numerical single-line CRD problem, not the legacy escape fit. The inner
boundary is explicit, rather than an artificial doubled optical-depth column.

For a normalized emission profile $\phi_i(\nu)$ and actual level gap $\Delta E_i$,
the adopted finite-width coefficients are
\begin{align}
 j_i(\nu)&=\frac{h\nu}{4\pi}n_{u,i}A_i\phi_i(\nu),\\
 \kappa_i(\nu)&=\frac{c^2 A_i\phi_i(\nu)}{8\pi\nu^2}
 \left[\frac{g_{u,i}}{g_{l,i}}n_{l,i}
 e^{(h\nu-\Delta E_i)/(kT)}-n_{u,i}\right].
\end{align}
The exponential is the thermodynamic profile correction: substituting
Boltzmann populations gives $j_i/\kappa_i=B_\nu$ at every frequency.
The profile center and level gap are supplied separately when a radiative
wavelength table has rounding differences from the level energies.
This enforces finite-width LTE, rather than relying on $B_\nu\simeq B_{\nu_i}$.
Positive opacity is required; coherent redistribution, moving-medium transfer,
and pressure wings are not supplied by this closure.

The cell's own trapped emission is eliminated algebraically from its SE
coefficients. If $\Lambda_{cc}(\nu)$ maps local emissivity to cell-average
$J_\nu$, its effective spontaneous loss is
$A_i\int\phi_i(1-\Lambda_{cc}\kappa_i)\,d\nu$.
The remaining field supplies excitation and stimulated de-excitation.
Incoming rays and other emitters are summed directly to avoid subtracting
two large trapped source functions. Photon balance is checked against the
actual face fluxes, including separately identified scattering removal.

Selected isolated OTS events are removed. Other species instead receive
the signed correction
\begin{equation}
 \Delta\Phi_\nu=\frac{4\pi\Delta\nu}{h\nu}
 (J_\nu^{\rm group}-J_\nu^{\rm continuum\ only}).
\end{equation}
Both solves use identical continuum coefficients and boundaries. Negative
entries represent screening of stellar photons; retaining only positive line
emission would miss that effect and double-count the baseline continuum.
Every species uses the same correction for photoionization, induced
recombination and heat. The line redistribution energy is the radiative
energy absorbed minus $\Delta E_i$ times the net excitation rate.
The optional scattering attenuation uses the same approximation as the H
continuum; it is not counted as true absorption. Continuum-grid consistency,
line-wing and spatial convergence remain separate validation requirements.

\section{An explicit infrared transfer alternative}
The optional \texttt{resolved\_infrared\_lines} selects a fixed list of
hydrogen shell transitions before solving. Intersecting fixed frequency
windows are grouped before constructing one radiation correction per group.
Each group uses the preceding formal-transfer and local-source
elimination machinery. The initial three-line experiment selects
$14\to10$, $14\to11$ and $15\to8$; the default list is empty. Selection never depends
on the current sign of an opacity: crossing zero must not change the model.
The original statistical-equilibrium equations remain, while selected rates
and isolated OTS events are replaced together. All continuum receivers use
the same signed field correction, once.

The optional \texttt{resolved\_infrared\_metal\_lines} supplies fixed
\texttt{(atom label, upper, lower)} entries with zero-based levels of the
retained atom. Their absorption and emission enter the same connected group
as hydrogen, and both atoms receive its radiative rates. The selected
metal lines' isolated OTS photons are removed together with their old rates;
the combined group supplies one continuum correction and energy ledger.
This also changes the metal line's isolated escape approximation to formal
transfer. Metal-only formal and joint H/metal controls therefore separate
that closure change from the additional overlap effect. No transition is
selected dynamically by its current population or opacity sign.

In this mode the scattering opacity is electron scattering only,
$\kappa_e=n_e\sigma_{\rm T}$. Let $\Lambda_c$ map cell emissivity to
cell-average intensity under extinction $\kappa_c+\kappa_e$, and let
$J_{b,c}$ be the boundary contribution. The broad continuum reservoir obeys
\begin{align}
 (I-\Lambda_c\kappa_e)J_{\rm bg}
     &=J_{b,c}+\Lambda_c\eta_c,\nonumber\\
 \eta_{\rm bg}&=\eta_c+\kappa_eJ_{\rm bg},\nonumber\\
 J_{\rm line}&=J_b(\kappa_{\rm tot})
       +\Lambda_{\rm tot}(\eta_{\rm bg}+\eta_{\rm line}),\qquad
 \kappa_{\rm tot}=\kappa_c+\kappa_e+\kappa_{\rm line}.
 \label{eq:infrared_background}
\end{align}
Products with an opacity act diagonally on the source-cell index. The
background linear solve has a checked residual. In the slab version, individual line opacity may
be negative only while total extinction stays positive and the replacement
rates remain nonnegative. This avoids a sign-dependent escape approximation;
it supplies no saturated-maser model. With Planck boundaries and LTE
coefficients, the background and full field recover the Planck function.

The reservoir returns continuum photons to the line band, but does not
follow frequency redistribution of the additional line photons after
electron scattering. Their signed removed energy is reported separately,
\begin{equation}
 Q_e=4\pi\int\kappa_e(J_{\rm line}-J_{\rm bg})\,d\nu,
 \label{eq:infrared_scattering_energy}
\end{equation}
and is not added as kinetic heat. A positive Green-operator norm bounds
successive omitted photon returns for the selected discrete cell/angle
operator. Its outside-band bound assumes the same electron opacity and
nonnegative other extinction; unknown inverted lines are outside that
statement. This conditional photon bound is not a bound on temperature,
population or spectral changes. Those require the complete atmosphere to
be re-solved and compared. With the legacy continuum option, the common field retains its earlier attenuation closure,
whereas the selected line groups use this scattering background and may use
finer angular quadrature. Consequently the group field need not equal the
common receiver field plus its correction: this is an explicit perturbative
alternative, not an exact shared-background claim. Matching and refining
both angular/background treatments, and checking a zero-line control and
band-width sensitivity, remain required. Rayleigh redistribution, electron
energy exchange and the adequacy of the fixed line list are also separate
obligations.

\section{Finite spherical paths for signed extinction}
The slab alternative above can fail even when amplification across the
planetary atmosphere is small: its grazing paths have unbounded length.
An actual eight-cell trial encounters negative total extinction in
$14\to11$ and is therefore outside that solver's domain. Raising a Newton
iteration limit cannot repair this geometric limitation. Population inversion
can arise from coupled hydrogen recombination and radiation; it is not by
itself proof of a numerical artifact \cite{zhu}. The following finite-shell
operator is an explicit alternative, whose atmospheric adequacy remains to
be tested.

Let $p$ be impact parameter and $r_{j-1},r_j$ the shell edges. The positive
half-chord length is
$\ell_j(p)=\sqrt{[r_j^2-p^2]_+}-\sqrt{[r_{j-1}^2-p^2]_+}$.
For constant extinction $\chi$ and emissivity $\eta$ on a segment,
\begin{align}
 \tau&=\chi\ell,\qquad
 \phi_1(\tau)=\frac{1-e^{-\tau}}{\tau},\qquad
 \phi_2(\tau)=\frac{1-\phi_1(\tau)}{\tau},\nonumber\\
 I_{\rm out}&=I_{\rm in}e^{-\tau}+\eta\ell\phi_1(\tau),\qquad
 \overline I=I_{\rm in}\phi_1(\tau)+\eta\ell\phi_2(\tau).
 \label{eq:spherical_segment}
\end{align}
The analytic limits $\phi_1(0)=1$ and $\phi_2(0)=1/2$ retain transparent
emission and derivatives. Signed $\chi$ is permitted without a floor;
overflow remains failure. An inward sweep starts at the outer boundary.
The outward sweep starts at the prescribed core intensity if the ray hits
the core, or at the inward turning-point intensity otherwise.

Writing $V_j=(r_j^3-r_{j-1}^3)/3$ for volume per solid angle, use the same
chords to compute shell-volume averages and face flux moments:
\begin{align}
 J_j&=\frac{1}{2V_j}\int_0^{r_j}p\,dp\,\ell_j(p)
          (\overline I_j^++\overline I_j^-),\nonumber\\
 H(r)&=\frac{1}{2r^2}\int_0^r p\,dp\,(I^+(r,p)-I^-(r,p)),\nonumber\\
 r_j^2H_j-r_{j-1}^2H_{j-1}&=V_j(\eta_j-\chi_jJ_j).
 \label{eq:spherical_average}
\end{align}
Split the impact integral at every shell edge. On a panel $[a,b]$, use
$p^2=b^2-(b^2-a^2)t^2$, with weight $p\,dp=(b^2-a^2)t\,dt$ and $0<t<1$,
to regularize tangency. The geometric-volume quadrature error is reported;
the exact $V_j$ is not renormalized to hide it.
The energy check sums absolute cell defects relative to the gross supplied
power at each frequency separately:
\begin{align}
 P_\nu&=\frac{r_0^2I_{\nu,\rm inner}+r_N^2I_{\nu,\rm outer}}{4}
       +\sum_j V_j\bigl[\eta_{\nu,j}+(-\chi_{\nu,j})_+J_{\nu,j}\bigr],\nonumber\\
 \epsilon_\nu&=\frac{\sum_j\left|
 r_j^2H_{\nu,j}-r_{j-1}^2H_{\nu,j-1}
 -V_j(\eta_{\nu,j}-\chi_{\nu,j}J_{\nu,j})\right|}{P_\nu}.
 \label{eq:spherical_energy_norm}
\end{align}
The zero-source, zero-field case has zero error. Finite stimulated gain is
included in supplied power. The maximum over frequencies and the geometric
volume error must each be below $10^{-9}$. This replaces a cell-local relative
ratio that became undefined for nearly extinguished, underflowing responses;
it does not guarantee relative accuracy of every faint field or population.
Angular and atmospheric convergence still require separate checks.
The implementation is
\path{escape.spherical_radiation.solve_spherical_transfer}.

Both boundaries are isotropic; the core absorbs incoming rays. A nonzero
collimated stellar beam is outside this spherical operator's scope. This
is static transfer at supplied populations, not a saturated-maser closure.
Selected line-group integration must use its actual Green operator: local
self-emission includes the far-side segment of the same shell. Reusing a
slab's diagonal response would be inconsistent. Matching this volume-mean
field to the common continuum, selecting overlapping line groups, and
re-solving the coupled populations and thermal state remain obligations.

\section{A shared infrared continuum control}
With \texttt{infrared\_continuum=spherical\_shell\_average}, a fixed band
uses the spherical coherent electron-scattering background for all retained
continuum receivers and line-pumping inputs. Its initial boundary is
$\lambda\geq3\,\mu{\rm m}$, independent of the selected line windows.
Both the band boundary and its lower-frequency limit require refinement.
The explicit \texttt{minimum\_continuum\_energy\_ev} sets that limit;
every selected pumping frequency must lie inside it. Endpoint responses
bracket the line centers without changing the integration weights.
The legacy continuum remains the default. This control changes the geometry;
it is not automatically equivalent to the substellar-column reference.

For prescribed radial coefficients, rotational averaging of a collimated
stellar field gives the same shell-volume mean as isotropic incoming
intensity $I_{\rm outer}=J_\star$. Indeed,
$\pi R^2 F_\star=4\pi^2R^2J_\star$ for $F_\star=4\pi J_\star$.
This identity does not prove equality of angularly varying populations or
temperature with their one-dimensional averages. The stellar beam is not
also added after this replacement.

Two right-hand sides of the continuum scattering solve separately give
the emission/inner-boundary field $E_\nu$ and the unit outer-boundary
response $R_\nu$. Thus
\begin{equation}
 J_{\nu,\rm bg}=E_\nu+R_\nu J_{\star,\nu}.
 \label{eq:infrared_shared_field}
\end{equation}
Only the smooth responses are interpolated to isolated line centers; the
stellar intensity is sampled at the exact center. Selected resolved groups
instead evaluate this same continuum closure on their own frequency nodes
and return $J_{\nu,\rm group}-J_{\nu,\rm bg}$. A zero-line control must
recover the common background, and spectral interpolation requires its own
test. This prevents arbitrary line-window width from defining a broad
continuum correction.

Hydrogen photoionization, induced recombination and thermal terms, and the
metal, helium, molecular and free--free receivers use the common field.
For ungrouped lines, its shell-volume mean is held constant at the existing
depth nodes before the original shielding/escape fits are applied. This
remains an angular and spatial reduction. Selected metal group members
participate reciprocally. Group photons are not added to ungrouped metal
bound--bound pumping without the corresponding opacity; remaining material
accidental H/metal infrared overlaps require expanded fixed selections.
Only electron scattering is included in this control: Rayleigh scattering
and redistribution need contribution bounds. The continuum solver is
\path{escape.spherical_radiation.solve_spherical_scattering}; its checked
energy ledger uses Eq.~\eqref{eq:spherical_energy_norm} with true absorption
and emission; coherent scattering is internal recycling, not supplied power
or heat. Shared routing and conservation do not establish atmospheric or
observational adequacy.

\begin{thebibliography}{9}
\raggedright
\bibitem{cloudy} Cloudy collaboration, C23.01 source distribution:
\texttt{rt\_escprob.cpp}, \texttt{rt\_continuum\_shield\_fcn.cpp},
\texttt{rt\_line\_one.cpp}, \texttt{rt\_ots.cpp},
\texttt{opacity\_add1subshell.cpp}, \texttt{rt\_continuum.cpp},
\texttt{two\_photon.cpp}. Altered-source notice and license are retained in
\path{src/escape/_licenses/Cloudy.txt}.
\bibitem{goldman} Goldman, S. P. \& Drake, G. W. F. (1981),
\emph{Physical Review A}, 24, 183.
\href{https://doi.org/10.1103/PhysRevA.24.183}{doi:10.1103/PhysRevA.24.183}.
\bibitem{lebedev} Lebedev, V. S. \& Beigman, I. L. (1998),
\emph{Physics of Highly Excited Atoms and Ions}, Springer, Eq.~8.30.
\href{https://doi.org/10.1007/978-3-642-72175-5}{doi:10.1007/978-3-642-72175-5}.
C23 implementation: \texttt{hydro\_vs\_rates.cpp}, \texttt{hydrocollid.cpp},
\texttt{helike\_cs.cpp}, \texttt{iso\_collide.cpp}.
\bibitem{netzer} Netzer, H., Elitzur, M. \& Ferland, G. J. (1985),
``Bowen fluorescence and He II lines in active galaxies and gaseous nebulae'',
\emph{Astrophysical Journal}, 299, 752--762, Eqs.~2--6 and Fig.~2.
\href{https://doi.org/10.1086/163741}{doi:10.1086/163741}.
\bibitem{zhu} Zhu, F.-Y., Wang, J., Qiu, Y., Zhu, Q.-F. \& Quan, D. (2025),
``An $nl$-model with full radiative transfer treatment for level populations
of hydrogen atoms in spherically symmetric H II region'',
\href{https://arxiv.org/abs/2509.12806}{arXiv:2509.12806}.
This provides physical context for coupled inversion and finite geometry;
its H II-region assumptions are not adopted as planetary validation.
\end{thebibliography}
\end{document}
